\section{Симметричные и кососимметричные формы}

\begin{definition}
    Пусть $V$ --- векторное пространство над полем $k$, и $B\colon V \times V \to k$ --- билинейная форма.

    \begin{enumerate}
        \item $B$ называется \emph{симметричной}, если $B(v, u) = B(u, v)$ для любых $v, u \in V$.
        \item $B$ называется \emph{кососимметричной}, если $B(v, v) = 0$ для любого $v \in V$.
        \item $B$ называется \emph{антисимметричной}, если $B(v, u) = -B(u, v)$ для любых $v, u \in V$.
    \end{enumerate}
\end{definition}


\begin{stmt}
    Антисимметричность влечёт кососимметричность. Обратное верно тогда и только тогда,
    когда $\mathrm{char}\, k \neq 2$.
    \begin{proof}
        \textbf{Антисимметричность $\Rightarrow$ кососимметричность:}
        так как $B(u+v, u+v) = 0$, раскрываем:
        \[
            0 = B(u,u) + B(u,v) + B(v,u) + B(v,v) = B(u,v) + B(v,u),
        \]
        откуда $B(u,v) = -B(v,u)$. Это верно при любой характеристике.

        \textbf{Кососимметричность $\Rightarrow$ антисимметричность} (при $\mathrm{char}\, k \neq 2$):
        подставляем $u = v$:
        \[
            B(v,v) = -B(v,v) \implies 2B(v,v) = 0.
        \]
        При $\mathrm{char}\, k \neq 2$ отсюда $B(v,v) = 0$.
        При $\mathrm{char}\, k = 2$ равенство $2B(v,v) = 0$ выполнено тождественно,
        поэтому $B(v,v)$ может быть ненулевым, и импликация в общем случае не верна.
    \end{proof}
\end{stmt}


\begin{remark}
    В поле характеристики $2$ выполняется $-1 = 1$, поэтому условие антисимметричности $B(v, u) = -B(u, v)$ совпадает с условием симметричности $B(v, u) = B(u, v)$. Таким образом, в характеристике $2$ антисимметричная форма --- это просто симметричная форма.
\end{remark}


\begin{stmt}
    ~\begin{enumerate}
        \item Если $B$ симметрична, $G_{\mathbf{v}}$ симметрична $G_{\mathbf{v}}^T = G_{\mathbf{v}}$: \[ B(v_i, v_j) = B(v_j, v_i). \]
        \item Если $B$ кососимметрична, $G_{\mathbf{v}}$ кососимметрична (и нули на диагонали)
            $G_{\mathbf{v}}^T = -G_{\mathbf{v}}$: \[ B(v_i, v_j) = -B(v_j, v_i). \]
        \item Если $B$ антисимметрична, то при $\mathrm{char}\, k = 2$ $G_{\mathbf{v}}$ симетрична 
            $G_{\mathbf{v}}^T = G_{\mathbf{v}}$: \[ B(v_i, v_j) = -B(v_j, v_i) = B(v_j, v_i), \]
            При $\mathrm{char}\, k \neq 2$ антисимметричность совпадает с кососимметричностью,
            и $G_{\mathbf{v}}^T = -G_{\mathbf{v}}$.

    \end{enumerate}
\end{stmt}



\section{Невырожденность}

\begin{definition}
    Билинейная форма $B$ называется \emph{невырожденной}, если $\forall\ u \neq 0\ \exists v\colon B(u,v) \neq 0$.
\end{definition}


\begin{theorem}
    Для билинейной формы $B$ (симметричной или кососимметричной) на конечномерном пространстве $V$ следующие условия эквивалентны:
    \begin{enumerate}
        \item $\forall\, u \neq 0\ \exists\, v\colon B(u, v) \neq 0$
        \item $\forall\, u \neq 0\ \exists\, v\colon B(v, u) \neq 0$
        \item Существует базис $\mathbf{e}$ с $\det G_{\mathbf{e}} \neq 0$
        \item Для любого базиса $\mathbf{e}$ выполняется $\det G_{\mathbf{e}} \neq 0$
        \item Левая корреляция --- изоморфизм
        \item Правая корреляция --- изоморфизм
    \end{enumerate}

    \begin{proof}
        TODO
        % Для симметричных и кососимметричных форм левая и правая корреляции отличаются знаком, поэтому $1 \Leftrightarrow 2$ и $5 \Leftrightarrow 6$.

        % \textbf{$3 \Rightarrow 4$.} Если $\mathbf{f}$ --- другой базис, то $G_{\mathbf{f}} = C_{\mathbf{e}\mathbf{f}}^T\, G_{\mathbf{e}}\, C_{\mathbf{e}\mathbf{f}}$, и $\det G_{\mathbf{f}} = (\det C_{\mathbf{e}\mathbf{f}})^2 \det G_{\mathbf{e}} \neq 0$.

        % \textbf{$4 \Rightarrow 1$.} Пусть $u \neq 0$. Дополним до базиса $\mathbf{e} = (u, e_2, \dots, e_n)$. Так как $\det G_{\mathbf{e}} \neq 0$, первая строка не нулевая, то есть $\exists\, j\colon B(u, e_j) \neq 0$.

        % \textbf{$2 \Rightarrow 5$.} Ядро левой корреляции $\hat{B}\colon V \to V^*$ равно $\{v \mid B(v, u) = 0\ \forall u\}$. По условию 2, $\ker\hat{B} = \{0\}$. Так как $\dim V = \dim V^*$, инъекция $\Rightarrow$ биекция.

        % \textbf{$5 \Rightarrow 3$.} Матрица левой корреляции в базисах $\mathbf{e}$ и $\mathbf{e}^*$ равна $G_{\mathbf{e}}^T$. Изоморфизм $\Rightarrow$ $\det G_{\mathbf{e}}^T \neq 0$ $\Rightarrow$ $\det G_{\mathbf{e}} \neq 0$.
    \end{proof}
\end{theorem}

% % ??
% \begin{remark}
%     Матрица левой корреляции в базисах $\mathbf{e}$ и $\mathbf{e}^*$ равна $G_{\mathbf{e}}^T$, а правой --- $G_{\mathbf{e}}$.

%     Для вычисления: $\hat{B}(e_i)$ --- функционал, который на $e_j$ даёт $B(e_j, e_i) = G_{ji}$. Разложив по двойственному базису, получаем $i$-й столбец $G^T$.
% \end{remark}



\begin{definition}    
    Ограничение $B|_U$ --- это билинейная форма на подпространстве $U \subseteq V$, определённая как
    \[
        B|_U \colon U \times U \to k, \quad B|_U(u, u') = B(u, u'),
    \]
    \[
        U \times U \hookrightarrow V \times V \xrightarrow{B} k.
    \]
\end{definition}


\begin{stmt}
    Ограничение $B$ (симметричная или кососимметричная форма) на дополнение к $\ker B$ невырождена.

    \begin{proof}
        $V = \ker B \oplus U$.
        Пусть $u \in U$, $u \neq 0$, и $B(u, v) = 0$ для всех $v \in U$. Покажем, что тогда $B(u, w) = 0$ для всех $w \in V$.

        Любой $w \in V$ представляется как $w = e + f$, где $e \in \ker B$, $f \in U$. Тогда
        \[
            B(u, w) = B(u, e) + B(u, f) = 0 + 0 = 0.
        \]
        Значит $u \in \ker B$. Но $u \in U$ и $\ker B \cap U = \{0\}$ (прямая сумма), откуда $u = 0$ --- противоречие.
    \end{proof}
\end{stmt}


\begin{stmt}
    Пусть $B$ --- симметричная или кососимметричная форма на $V$. 
    Тогда $B$ задаёт невырожденную (косо)симметричную форму $\bar{B}$ на $V / \ker B$:
    \[
        \bar{B}([v], [u]) \coloneqq B(v, u).
    \]

    \begin{proof}
        \textbf{Корректность.} Пусть $v' \sim v$, $u' \sim u$, то есть $v' - v, \, u' - u \in \ker B$. Тогда:
        \[
            B(v', u') = B(v + \underbrace{(v' - v)}_{\in\, \ker B},\, u + \underbrace{(u' - u)}_{\in\, \ker B}) = B(v, u),
        \]
        так как все слагаемые с элементами ядра обращаются в ноль.

        \textbf{Невырожденность.} Пусть $[v] \neq 0$, то есть $v \notin \ker B$.
        Тогда $\exists\, u \in V\colon B(v, u) \neq 0$, и значит $\bar{B}([v], [u]) \neq 0$.

        \textbf{Симметричность / кососимметричность} наследуется, поскольку $\bar{B}$ определяется через $B$.
    \end{proof}
\end{stmt}


\begin{definition}
    Пусть $U \subseteq V$ --- подпространство, $B$ --- симметричная или кососимметричная форма. \emph{Ортогональное дополнение} $U$ относительно $B$:
    \[
        U^\perp \coloneqq \{v \in V \mid B(v, u) = 0\ \forall\, u \in U\}.
    \]
\end{definition}

\begin{theorem}
    Если $B\big|_U$ невырожденна, то $V = U \oplus U^\perp$.
    \begin{proofsteps}
        \item Прямая сумма: ($U \cap U^\perp = \{0\}$): если $v \in U \cap U^\perp$, то $B(v, u) = 0$ для всех $u \in U$. Так как $v \in U$, это означает, что $v \in \ker(B\big|_U) = \{0\}$.

        \item $V = U + U^\perp$. Так как $B\big|_U$ невырожденна, корреляция $\hat{B}_U\colon U \to U^*$, $u \mapsto B(u, \cdot)\big|_U$ --- изоморфизм.

        Для произвольного $v \in V$ рассмотрим функционал $\varphi_v\colon U \to k$, $\varphi_v(u) = B(v, u)$. По изоморфизму $\exists\, p_U(v) \in U$ такой, что
        \[
            B(p_U(v), u) = B(v, u) \quad \forall\, u \in U.
        \]
        Тогда $v - p_U(v) \in U^\perp$, поскольку
        \[
            B(v - p_U(v),\, u) = B(v, u) - B(p_U(v), u) = 0 \quad \forall\, u \in U.
        \]
        Следовательно, $v = \underbrace{p_U(v)}_{\in\, U} + \underbrace{(v - p_U(v))}_{\in\, U^\perp}$.
    \end{proofsteps}
\end{theorem}


\begin{remark}
    В отличие от евклидова случая, если $B\big|_U$ вырождена, ортогональное дополнение $U^\perp$ может пересекаться с $U$ нетривиально, и $U + U^\perp \neq V$ в общем случае.
\end{remark}

